% GENERATED FILE. Edit canonical lessons, then run npm run generate:books.
% canonical-source-hash: fnv1a64:8e4aa997f6b23d20
% canonical-lessons: RU-C404-kapot, RU-C404-dvigatel, RU-C404-akkumulyator, RU-C404-zavodit, RU-C404-glokhnut

\chapter{Under the bonnet}
\label{ch:ru-under-the-bonnet}
\hlchaptermodality{404}

\begin{chapteropening}
\textbf{By the end of this chapter:} \emph{I can talk about what is under the bonnet: the engine, the battery, starting the car and an engine that stalls.}

Complete the last lesson of chapter 404: I can talk about what is under the bonnet: the engine, the battery, starting the car and an engine that stalls.
\end{chapteropening}

\section[\texorpdfstring{\ru{капот} (kapót)}{капот (kapót)}]{\ru{капот} (kapót) — a bonnet (of a car); a hood}

\label{lesson:RU-C404-kapot}



\begin{quote}
Before the new one: say the Russian for to speed up, then the Russian for to sound the horn.
\end{quote}

\subsection*{You'll want to know: \ru{капот}}
\textbf{\ru{капот}} (\emph{kapót}) — \textquotedblleft{}a bonnet (of a car); a hood\textquotedblright{}.

\begin{quote}
\ru{Водитель} \ru{вышел} \ru{и} \ru{открыл} \ru{капот}.
\end{quote}

\begin{quote}
\emph{Vodítel výshel i otkrýl kapót.}
\end{quote}

\begin{quote}
The driver got out and opened the bonnet.
\end{quote}

Masculine. In American English, the hood.

\begin{cousinweb}
\emph{\ru{капот}} was borrowed from French \emph{capote}, a hooded cloak (EtymDB): the engine wears a hood. EtymDB traces \emph{capote} back to Latin \emph{caput} 'head'.
\end{cousinweb}

\begin{grammarlens}[title={What you've built}]
One more word for telling, explaining and describing at B1.
\end{grammarlens}

\subsection*{Your turn}
\begin{itemize}
\raggedright
  \item \emph{Say it:} \emph{kapót}
  \item \emph{Say it:} \emph{kapót}, once more
  \item \emph{Say it:} a sentence of your own with \emph{kapót}
  \item \emph{Recall:} say the Russian for to speed up, then the Russian for to sound the horn, then say \emph{kapót} again
\end{itemize}

\begin{tcolorbox}[breakable,colback=teal!4,colframe=teal!35!black,title={Before you move on}]
What does \ru{капот} mean? (\textquotedblleft{}A bonnet (of a car); a hood\textquotedblright{}.) Say it once more.
\end{tcolorbox}
\section[\texorpdfstring{\ru{двигатель} (dvígatel)}{двигатель (dvígatel)}]{\ru{двигатель} (dvígatel) — an engine, a motor}

\label{lesson:RU-C404-dvigatel}



\begin{quote}
Before the new one: say the Russian for to sound the horn, then the Russian for a bonnet.
\end{quote}

\subsection*{You'll want to know: \ru{двигатель}}
\textbf{\ru{двигатель}} (\emph{dvígatel}) — \textquotedblleft{}an engine, a motor\textquotedblright{}.

\begin{quote}
\ru{У} \ru{новой} \ru{машины} \ru{тихий} \ru{двигатель}.
\end{quote}

\begin{quote}
\emph{U nóvoy mashíny tíkhiy dvígatel.}
\end{quote}

\begin{quote}
The new car has a quiet engine.
\end{quote}

Masculine, though it ends in -\ru{ь}: \emph{\ru{мой} \ru{двигатель}}.

\begin{cousinweb}
\emph{\ru{двигатель}} is built on \emph{\ru{двигать}} 'to move' with the ending \emph{-\ru{тель}}, 'the thing that does it'. \emph{\ru{Двигать}} is inherited from Proto-Slavic \emph{dvigati} 'to lift, to move' (EtymDB).
\end{cousinweb}

\begin{grammarlens}[title={What you've built}]
One more word for telling, explaining and describing at B1.
\end{grammarlens}

\subsection*{Your turn}
\begin{itemize}
\raggedright
  \item \emph{Say it:} \emph{dvígatel}
  \item \emph{Say it:} \emph{dvígatel}, once more
  \item \emph{Say it:} a sentence of your own with \emph{dvígatel}
  \item \emph{Recall:} say the Russian for to sound the horn, then the Russian for a bonnet, then say \emph{dvígatel} again
\end{itemize}

\begin{tcolorbox}[breakable,colback=teal!4,colframe=teal!35!black,title={Before you move on}]
What does \ru{двигатель} mean? (\textquotedblleft{}An engine, a motor\textquotedblright{}.) Say it once more.
\end{tcolorbox}
\section[\texorpdfstring{\ru{аккумулятор} (akkumulyátor)}{аккумулятор (akkumulyátor)}]{\ru{аккумулятор} (akkumulyátor) — a car battery; a rechargeable battery}

\label{lesson:RU-C404-akkumulyator}



\begin{quote}
Before the new one: say the Russian for a bonnet, then the Russian for an engine.
\end{quote}

\subsection*{You'll want to know: \ru{аккумулятор}}
\textbf{\ru{аккумулятор}} (\emph{akkumulyátor}) — \textquotedblleft{}a car battery; a rechargeable battery\textquotedblright{}.

\begin{quote}
\ru{Ночью} \ru{был} \ru{сильный} \ru{холод}, \ru{и} \ru{аккумулятор} \ru{сел}.
\end{quote}

\begin{quote}
\emph{Nóchyu byl sílnyy khólod, i akkumulyátor sel.}
\end{quote}

\begin{quote}
It was very cold in the night, and the battery went flat.
\end{quote}

\emph{\ru{Аккумулятор} \ru{сел}} — the battery is flat. Small batteries for a remote are \emph{\ru{батарейки}}.

\begin{cousinweb}
\emph{\ru{аккумулятор}} is usually traced to Latin \emph{accumulātor} 'one who heaps up', from \emph{accumulāre} 'to pile up' (EtymDB lists the Latin words): a battery stores up charge. EtymDB does not list the Russian word, and the route into Russian, probably through German, is not settled.
\end{cousinweb}

\begin{grammarlens}[title={What you've built}]
One more word for telling, explaining and describing at B1.
\end{grammarlens}

\subsection*{Your turn}
\begin{itemize}
\raggedright
  \item \emph{Say it:} \emph{akkumulyátor}
  \item \emph{Say it:} \emph{akkumulyátor}, once more
  \item \emph{Say it:} a sentence of your own with \emph{akkumulyátor}
  \item \emph{Recall:} say the Russian for a bonnet, then the Russian for an engine, then say \emph{akkumulyátor} again
\end{itemize}

\begin{tcolorbox}[breakable,colback=teal!4,colframe=teal!35!black,title={Before you move on}]
What does \ru{аккумулятор} mean? (\textquotedblleft{}A car battery; a rechargeable battery\textquotedblright{}.) Say it once more.
\end{tcolorbox}
\section[\texorpdfstring{\ru{заводить} (zavodít)}{заводить (zavodít)}]{\ru{заводить} (zavodít) — to start (an engine); to wind up (a clock)}

\label{lesson:RU-C404-zavodit}



\begin{quote}
Before the new one: say the Russian for an engine, then the Russian for a car battery.
\end{quote}

\subsection*{You'll want to know: \ru{заводить}}
\textbf{\ru{заводить}} (\emph{zavodít}) — \textquotedblleft{}to start (an engine); to wind up (a clock)\textquotedblright{}.

\begin{quote}
\ru{Каждое} \ru{утро} \ru{сосед} \ru{заводит} \ru{машину} \ru{в} \ru{семь} \ru{часов}.
\end{quote}

\begin{quote}
\emph{Kázhdoye útro soséd zavódit mashínu v sem chasóv.}
\end{quote}

\begin{quote}
Every morning the neighbour starts his car at seven.
\end{quote}

Imperfective; the perfective is \emph{\ru{завести}} (\emph{\ru{заведу}, \ru{заведёшь}}). Present \emph{\ru{завожу}, \ru{заводишь}}. \emph{\ru{Машина} \ru{не} \ru{заводится}} — the car will not start.

\begin{cousinweb}
\emph{\ru{заводить}} is built from \emph{\ru{за}-} and \emph{\ru{водить}} 'to lead'. \emph{\ru{Водить}} is inherited from Proto-Slavic \emph{voditi} 'to lead' (EtymDB).
\end{cousinweb}

\begin{grammarlens}[title={What you've built}]
One more word for telling, explaining and describing at B1.
\end{grammarlens}

\subsection*{Your turn}
\begin{itemize}
\raggedright
  \item \emph{Say it:} \emph{zavodít}
  \item \emph{Say it:} \emph{zavodít}, once more
  \item \emph{Say it:} a sentence of your own with \emph{zavodít}
  \item \emph{Recall:} say the Russian for an engine, then the Russian for a car battery, then say \emph{zavodít} again
\end{itemize}

\begin{tcolorbox}[breakable,colback=teal!4,colframe=teal!35!black,title={Before you move on}]
What does \ru{заводить} mean? (\textquotedblleft{}To start (an engine); to wind up (a clock)\textquotedblright{}.) Say it once more.
\end{tcolorbox}
\section[\texorpdfstring{\ru{глохнуть} (glókhnut)}{глохнуть (glókhnut)}]{\ru{глохнуть} (glókhnut) — to stall (of an engine); to go deaf}

\label{lesson:RU-C404-glokhnut}



\begin{quote}
Before the new one: say the Russian for a car battery, then the Russian for to start.
\end{quote}

\subsection*{You'll want to know: \ru{глохнуть}}
\textbf{\ru{глохнуть}} (\emph{glókhnut}) — \textquotedblleft{}to stall (of an engine); to go deaf\textquotedblright{}.

\begin{quote}
\ru{Старый} \ru{мотор} \ru{часто} \ru{глохнет} \ru{на} \ru{светофоре}.
\end{quote}

\begin{quote}
\emph{Stáryy motór chásto glókhnet na svetofóre.}
\end{quote}

\begin{quote}
The old engine often stalls at the traffic lights.
\end{quote}

Imperfective; perfective \emph{\ru{заглохнуть}}. Past \emph{\ru{глох}, \ru{глохла}}: the -\ru{ну}- drops out.

\begin{cousinweb}
\emph{\ru{глохнуть}} is inherited from Proto-Slavic \emph{glŭxnǫti} (EtymDB), the verb that goes with \emph{\ru{глухой}} 'deaf': an engine that stalls goes silent.
\end{cousinweb}

\begin{grammarlens}[title={What you've built}]
One more word for telling, explaining and describing at B1.
\end{grammarlens}

\subsection*{Your turn}
\begin{itemize}
\raggedright
  \item \emph{Say it:} \emph{glókhnut}
  \item \emph{Say it:} \emph{glókhnut}, once more
  \item \emph{Say it:} a sentence of your own with \emph{glókhnut}
  \item \emph{Recall:} say the Russian for a car battery, then the Russian for to start, then say \emph{glókhnut} again
\end{itemize}

\begin{tcolorbox}[breakable,colback=teal!4,colframe=teal!35!black,title={Before you move on}]
What does \ru{глохнуть} mean? (\textquotedblleft{}To stall (of an engine); to go deaf\textquotedblright{}.) Say it once more.
\end{tcolorbox}
